% Week 8: Ethical Concerns (AI, Plagiarism) and Review of KS Protocol
% BA-Eindwerkstuk Seminar - Leiden University
% Dr. Steven Denney

\documentclass[aspectratio=169,11pt]{beamer}

% ============================================
% Packages
% ============================================
\usepackage{tikz}
\usepackage{graphicx}
\usepackage{hyperref}
\usepackage{fontspec}
\usepackage{booktabs}

% Load custom theme
\usepackage{beamerthemethesis}

% ============================================
% Metadata
% ============================================
\title{BA-Eindwerkstuk Seminar}
\subtitle{Week 8: Ethical Concerns and Review of the Thesis Protocol}
\author{Dr. Steven Denney}
\institute{Korean Studies \\ Leiden University}
\date{April 17, 2026}

% ============================================
% Document
% ============================================
\begin{document}

% --------------------------------------------
% Title Slide
% --------------------------------------------
\begin{frame}[plain,noframenumbering]
    \titlepage
\end{frame}

% ============================================
\section{Data Check-In}
% ============================================

\begin{frame}{Today's Agenda}
    \begin{enumerate}
        \item Data check-in: revisiting FAIR principles
        \item AI and ethics in research
        \item Discussion: your experience with AI tools
        \item Review of the Thesis Protocol and Assignment \#3
    \end{enumerate}
    \vspace{0.5cm}
    \begin{block}{Context}
        You have submitted Assignment \#2. You are deep in your research. Today we pause to address two things: the ethics of how you work, and the expectations for what comes next.
    \end{block}
\end{frame}

\begin{frame}{Where Are You With Your Data?}
    Assignment \#2 asked you to present your methodology and analytical framework. Now comes the harder part: \textbf{doing the analysis}.
    \vspace{0.5cm}

    \textbf{Honest self-assessment:}
    \begin{itemize}
        \item Do you actually \textit{have} your data---or are you still planning to collect it?
        \item Is it organized in a way that lets you work with it?
        \item Could someone else understand what you collected and why?
        \item Are you confident you can analyze it with the methods you described?
    \end{itemize}
    \vspace{0.3cm}
    \begin{alertblock}{No judgment}
        If the answer to any of these is ``no,'' that is fine---but you need to address it \textit{now}, not in Week 10.
    \end{alertblock}
\end{frame}

\begin{frame}{Revisiting FAIR Principles}
    In Week 2 we introduced FAIR as a way to think structurally about data. Now that you are \textit{in} your research, let's check in.
    \vspace{0.3cm}
    \begin{table}
        \small
        \begin{tabular}{lp{0.65\textwidth}}
            \toprule
            \textbf{Principle} & \textbf{Ask Yourself} \\
            \midrule
            \textbf{Findable} & Can you locate all your sources? Are files named clearly? Could you find a specific document in under a minute? \\
            \addlinespace
            \textbf{Accessible} & Is your data backed up? Could you recover everything if your laptop died tomorrow? \\
            \addlinespace
            \textbf{Interoperable} & Are your files in standard formats? Could your supervisor open and read them? \\
            \addlinespace
            \textbf{Reusable} & Have you documented \textit{where} each source came from, \textit{when} you collected it, and \textit{how}? \\
            \bottomrule
        \end{tabular}
    \end{table}
\end{frame}

\begin{frame}{Exercise: Data Readiness Check}
    \textbf{Take 5 minutes. Write brief answers to:}
    \vspace{0.3cm}
    \begin{enumerate}
        \item What data do you have \textit{in hand} right now?
        \item What data do you still need to collect?
        \item Where is it stored and how is it organized?
        \item What is your plan for the next three weeks to get from where you are to a draft empirical chapter?
    \end{enumerate}
    \vspace{0.5cm}
    \begin{block}{Discussion}
        Let's hear from a few of you. What is going well? Where are you stuck?
    \end{block}
\end{frame}

\begin{frame}{Data Collection and Transparency}
    Your analytical framework describes \textit{what} you will do. But your reader also needs to know:
    \vspace{0.3cm}
    \begin{itemize}
        \item \textbf{How} you collected or generated your data
        \item \textbf{Why} you selected these particular sources (and not others)
        \item \textbf{What limitations} your data has---and how those affect your findings
    \end{itemize}
    \vspace{0.3cm}
    This is not just methodological rigor. It is an \textbf{ethical obligation}: transparency about your evidence is foundational to academic integrity.
    \vspace{0.3cm}
    \begin{exampleblock}{Remember}
        Open science does not mean you must share everything publicly. It means you must be \textit{honest} about what you did and how you did it.
    \end{exampleblock}
\end{frame}

% ============================================
\section{AI and Ethics in Research}
% ============================================

\begin{frame}{Why Talk About This Now?}
    You are in the middle of independent research. You are writing, reading, analyzing, and making decisions every day.
    \vspace{0.5cm}

    Generative AI tools are widely available. Many of you may already be using them---for brainstorming, for writing, for translation, or for other tasks.
    \vspace{0.5cm}

    \textbf{This is not a scare talk.} It is a practical conversation about:
    \begin{itemize}
        \item What the university policy actually says
        \item What responsible use looks like
        \item Where the lines are
    \end{itemize}
\end{frame}

\begin{frame}{The Faculty Policy: Five Basic Principles}
    The Faculty of Humanities \textit{Guidelines for the Use of GenAI in Assessment} (August 2025) rest on five principles:
    \vspace{0.3cm}
    \begin{enumerate}
        \item \textbf{Integrity} --- Your degree represents specific academic skills. The learning process matters more than the product.
        \item \textbf{Transparency} --- You must always be open about your working methods.
        \item \textbf{Independence} --- You must demonstrate that \textit{you} achieved the learning objectives.
        \item \textbf{Fraud} --- Any action that hinders correct assessment of your knowledge is fraud.
        \item \textbf{Responsibility} --- You have full responsibility for what you submit.
    \end{enumerate}
\end{frame}

\begin{frame}{The Default Rule}
    \begin{alertblock}{The default position is prohibition}
        The use of GenAI products or content in \textbf{any form of assessment} is \textbf{not permitted}, unless your teacher has \textbf{explicitly} indicated which uses are allowed and under what conditions.
    \end{alertblock}
    \vspace{0.5cm}
    Key implications:
    \begin{itemize}
        \item Permission in one course does \textbf{not} carry over to another
        \item If your teacher has not said it is allowed, assume it is \textbf{not}
        \item When in doubt: \textbf{ask your supervisor}
    \end{itemize}
\end{frame}

\begin{frame}{What You CAN Do}
    The following are \textbf{generally permitted}, even without explicit instructor approval:
    \vspace{0.3cm}
    \begin{itemize}
        \item \textbf{Spelling and grammar checks} --- Tools like Grammarly or built-in spell checkers are fine, unless language skills are being assessed
        \item \textbf{Machine translation of your own text} --- e.g., translating a sentence you wrote from Dutch to English using DeepL
    \end{itemize}
    \vspace{0.3cm}
    With \textbf{explicit supervisor permission} and \textbf{full disclosure}:
    \begin{itemize}
        \item Using AI to brainstorm or generate ideas (not final text)
        \item Using AI to help understand a concept (as a learning aid)
        \item Using AI to assist with coding or data formatting tasks
    \end{itemize}
    \vspace{0.3cm}
    \textit{In all cases: disclose, disclose, disclose.}
\end{frame}

\begin{frame}{Best Practices}
    \begin{enumerate}
        \item \textbf{Verify everything} --- AI tools produce plausible-sounding text that may be factually wrong. LLMs generate \textit{language}, not \textit{knowledge}. Always check.
        \item \textbf{Be aware of hallucinations} --- AI tools fabricate citations, invent sources, and present false claims confidently. \textit{Never} trust an AI-generated reference without verifying it exists.
        \item \textbf{Watch for bias} --- AI output can contain and reinforce stereotypes (ethnic, gender, cultural). Be critical.
        \item \textbf{Keep records} --- Maintain a log of all prompts and outputs. Faculty policy requires you to hand these over on request.
        \item \textbf{Disclose all use} --- In your thesis, describe any AI tools used and how.
        \item \textbf{Treat AI as a tool, not an author} --- Your thesis must reflect \textit{your} thinking, \textit{your} analysis, and \textit{your} voice.
    \end{enumerate}
\end{frame}

\begin{frame}{What You CANNOT Do}
    \begin{alertblock}{Prohibited uses}
        \begin{itemize}
            \item \textbf{Submit AI-generated text as your own work} --- including copying, paraphrasing, or translating AI output without citation
            \item \textbf{Upload confidential or copyrighted material} to AI tools --- this includes interview transcripts, unpublished data, and others' work
            \item \textbf{Use AI to rewrite someone else's text} --- even with paraphrasing, this is fraud
            \item \textbf{Reproduce AI-generated material without citation} --- if you use it, cite it
            \item \textbf{Assume blanket permission} --- approval in one course or for one task does not extend to others
        \end{itemize}
    \end{alertblock}
\end{frame}

\begin{frame}{When AI Use Is Permitted: How to Cite}
    If your supervisor explicitly allows a specific use of GenAI, you must cite it. The faculty provides citation formats:
    \vspace{0.3cm}
    \footnotesize
    \begin{table}
        \begin{tabular}{p{0.12\textwidth}p{0.55\textwidth}p{0.2\textwidth}}
            \toprule
            \textbf{Tool} & \textbf{Reference} & \textbf{In-text} \\
            \midrule
            ChatGPT & OpenAI. (2023). ChatGPT (Sept 25 version) [Large Language Model]. \url{https://chat.openai.com/} & (OpenAI, 2023) \\
            \addlinespace
            Gemini & Google AI. (2023). Gemini (Oct 23 version) [Large language model]. Google. \url{https://gemini.google.com/} & (Google AI, 2023) \\
            \bottomrule
        \end{tabular}
    \end{table}
    \normalsize
    \vspace{0.3cm}
    \begin{block}{Important}
        You must also maintain a record of all prompts and AI output, to be provided to your supervisor on request.
    \end{block}
\end{frame}

\begin{frame}{Machine Translation: A Common Case}
    Machine translation (DeepL, Google Translate) is \textbf{generally allowed} for translating text you wrote yourself.
    \vspace{0.3cm}

    \textbf{How to cite it:}
    \vspace{0.2cm}

    \begin{exampleblock}{Example (from faculty guidelines)}
        \footnotesize
        ``O fenomeno da traducao indireta\ldots'' [The phenomenon of indirect translation is a manifestation of the marginality of a language\ldots] (Pieta 2013, p.~40; our translation supported by DeepL).
    \end{exampleblock}
    \vspace{0.3cm}

    \textbf{Not allowed:}
    \begin{itemize}
        \item Using machine translation to translate \textit{someone else's text} and presenting it as your own work
        \item Using translation tools when translation skill is being assessed
    \end{itemize}
\end{frame}

\begin{frame}{Plagiarism: A Brief Reminder}
    Plagiarism is presenting another person's work or ideas as your own. This includes:
    \vspace{0.3cm}
    \begin{itemize}
        \item Using direct quotations without citation
        \item Paraphrasing without acknowledgment
        \item Appropriating someone else's arguments or structure
        \item Failing to use quotation marks for direct quotes
        \item Internet and digital sources are held to the \textbf{same standard} as print sources
    \end{itemize}
    \vspace{0.3cm}
    Plagiarism results in course failure. See the \textcolor{SlateBlue}{\href{https://www.organisatiegids.universiteitleiden.nl/en/regulations/general/plagiarism}{Regulations on Plagiarism}} for the full policy.
    \vspace{0.3cm}

    \begin{block}{The connection}
        The AI guidelines and the plagiarism regulations share the same foundation: \textbf{your work must be your own}, and you must be \textbf{transparent} about how it was produced.
    \end{block}
\end{frame}

% ============================================
\section{Discussion}
% ============================================

\begin{frame}{Scenarios: What Would You Do?}
    \textbf{Discuss in pairs, then we'll talk as a group.}
    \vspace{0.3cm}
    \begin{enumerate}
        \item You use ChatGPT to brainstorm three possible ways to organize your findings chapter. You don't copy any text---you just use the structure as inspiration. Do you need to disclose this?
        \vspace{0.2cm}
        \item You paste a paragraph from a Korean newspaper article into DeepL and use the English translation in your thesis, with a citation for the original article. Is this acceptable?
        \vspace{0.2cm}
        \item You ask an AI tool to ``improve the academic tone'' of a paragraph you wrote. It rewrites several sentences substantially. You use the rewritten version. What are the issues?
        \vspace{0.2cm}
        \item You upload your interview transcripts to an AI tool to help identify themes. What concerns should you have?
    \end{enumerate}
\end{frame}

\begin{frame}{Open Conversation}
    \textbf{Let's talk honestly:}
    \vspace{0.5cm}
    \begin{itemize}
        \item Have you used AI tools in your research process so far?
        \item What for? What worked? What didn't?
        \item Where do you feel uncertain about the boundaries?
        \item What would help you make better decisions about AI use going forward?
    \end{itemize}
    \vspace{0.5cm}
    \begin{block}{Ground rule}
        This is a learning conversation, not an investigation. The goal is to leave this room with a clear understanding of the policy and confidence in how to apply it.
    \end{block}
\end{frame}

\begin{frame}[plain,noframenumbering]
    \vspace{2.5cm}
    \begin{center}
        {\color{AccentBlue}\Huge\bfseries Break}
        \vspace{0.3cm}

        {\color{DustyRose}\rule{3cm}{0.8pt}}
        \vspace{0.5cm}

        {\color{TextMuted}\large 10 minutes}
    \end{center}
\end{frame}

% ============================================
\section{Thesis Protocol Review}
% ============================================

\begin{frame}{Where We Are}
    \begin{table}
        \small
        \begin{tabular}{lll}
            \toprule
            \textbf{Milestone} & \textbf{Date} & \textbf{Status} \\
            \midrule
            Boot camp (Weeks 1--4) & Feb. 06 -- Feb. 27 & \textcolor{SageGreen}{Complete} \\
            Assignment \#1: Revised Proposal & Mar. 13 & \textcolor{SageGreen}{Complete} \\
            Peer review (Week 7) & Apr. 10 & \textcolor{SageGreen}{Complete} \\
            Assignment \#2: Preliminary Draft & Apr. 03 & \textcolor{SageGreen}{Complete} \\
            \textbf{You are here} & \textbf{Apr. 17} & --- \\
            Assignment \#3: Empirical Draft & May 06 & \textcolor{Terracotta}{In 3 weeks} \\
            Final Manuscript & Jun. 01 & \textcolor{Terracotta}{In 7 weeks} \\
            \bottomrule
        \end{tabular}
    \end{table}
    \vspace{0.3cm}
    \begin{alertblock}{The trajectory}
        You have written a proposal and a preliminary draft. The next step is to show that you can \textbf{do the research}---collect data, analyze it, and present findings.
    \end{alertblock}
\end{frame}

\begin{frame}{Assignment \#3: Empirical Draft}
    \textbf{Due: May 06, 2026} \hfill \textit{5,000+ words}
    \vspace{0.3cm}

    This is a draft of your empirical chapter. It should:
    \begin{itemize}
        \item Apply the research design laid out in earlier assignments
        \item Present the initial stages of data analysis
        \item Demonstrate that data aligns with your question and methodology
        \item Show preliminary findings with emphasis on how evidence addresses the research problem
        \item Build upon and integrate with your literature review and analytical framework
    \end{itemize}
    \vspace{0.3cm}
    \begin{block}{Key point}
        The empirical chapter must be \textbf{coherently linked} to your prior sections. It is not a standalone piece---it shows a logical connection between your design and your emerging analysis.
    \end{block}
\end{frame}

\begin{frame}{Assignment \#3: How It Is Assessed}
    \footnotesize
    \begin{table}
        \begin{tabular}{p{0.2\textwidth}p{0.7\textwidth}}
            \toprule
            \textbf{Criterion} & \textbf{What Distinguishes Strong Work} \\
            \midrule
            Empirical Analysis & Analysis is aligned with the research question and shows insight into the data; well-supported claims \\
            \addlinespace
            Use of Evidence & Appropriate primary data; well-selected excerpts, quotes, or figures that fully support the argument \\
            \addlinespace
            Integration with Literature \& Method & Excellent integration with prior sections; empirical chapter builds on and completes the earlier design \\
            \addlinespace
            Structure \& Clarity & Clear organization, strong transitions, academic tone and style \\
            \bottomrule
        \end{tabular}
    \end{table}
    \normalsize
    \vspace{0.3cm}
    The full rubric is available on the \textcolor{SlateBlue}{\href{https://scdenney.github.io/baks_thesis-seminar/assignments/}{course website}}.
\end{frame}

\begin{frame}{The Four Assessment Criteria (Protocol)}
    Your final thesis is assessed on four criteria. These apply to Assignment \#3 as well:
    \vspace{0.3cm}
    \begin{enumerate}
        \item \textbf{Knowledge and insight} --- Research question, motivation, literature review. Does the question reflect insight into key discussions in the field?
        \vspace{0.15cm}
        \item \textbf{Application of knowledge} --- Data, methods, analysis. Is there critical analysis of primary sources? Are methods effective and justified?
        \vspace{0.15cm}
        \item \textbf{Reaching conclusions} --- Synthesis, contribution, limitations. Are conclusions logical and well-founded? Is empirical analysis central?
        \vspace{0.15cm}
        \item \textbf{Communication} --- Writing, structure, citations. Is the language competent? Is the structure clear? Are citations correct?
    \end{enumerate}
\end{frame}

\begin{frame}{Example Thesis Structure}
    \begin{columns}[T]
        \begin{column}{0.48\textwidth}
            \textbf{I. Introduction}
            \begin{itemize}
                \footnotesize
                \item Research question
                \item Problem, gap, motivation
                \item Proposed research plan
                \item Thesis roadmap
            \end{itemize}
            \vspace{0.2cm}
            \textbf{II. Literature Review}
            \begin{itemize}
                \footnotesize
                \item LR section I, II, \ldots \textit{n}
            \end{itemize}
            \vspace{0.2cm}
            \textbf{III. Analytical Framework}
            \begin{itemize}
                \footnotesize
                \item Design, data, and methods
                \item Methodological concerns and scope
            \end{itemize}
        \end{column}
        \begin{column}{0.48\textwidth}
            \textbf{IV. Empirical Findings} \textcolor{Terracotta}{$\leftarrow$ Assignment \#3}
            \begin{itemize}
                \footnotesize
                \item Outcome I, II, \ldots \textit{n}
            \end{itemize}
            \vspace{0.2cm}
            \textbf{V. Conclusion and Discussion}
            \begin{itemize}
                \footnotesize
                \item Summary of research and findings
                \item Research contributions
                \item Limitations
                \item Future research
            \end{itemize}
            \vspace{0.2cm}
            \textbf{VI. Bibliography}
            \vspace{0.2cm}

            \textbf{VII. Appendices}
        \end{column}
    \end{columns}
\end{frame}

\begin{frame}{Grade Descriptors: What Do the Numbers Mean?}
    \begin{table}
        \small
        \begin{tabular}{lp{0.7\textwidth}}
            \toprule
            \textbf{Grade} & \textbf{What It Means} \\
            \midrule
            9--10 & Outstanding; original, independent thinking; rigorous argument supported by wide evidence \\
            \addlinespace
            8--8.9 & Excellent understanding; independent thought; strong, well-organized argument \\
            \addlinespace
            7--7.9 & Good to very good; most criteria met \\
            \addlinespace
            6--6.9 & Satisfactory; reasonable argument with some shortcomings but no fundamental errors \\
            \addlinespace
            5.1--5.9 & \textbf{The faculty does not issue grades in this range} \\
            \addlinespace
            3--5.0 & Inadequate; substantial omissions \\
            \bottomrule
        \end{tabular}
    \end{table}
    \vspace{0.2cm}
    \textit{To pass the seminar: overall mark of 5.50 (= 6) or higher, with Assignment \#3 at 5.5+.}
\end{frame}

% ============================================
\section{Looking Ahead}
% ============================================

\begin{frame}{The Path Forward}
    \begin{table}
        \small
        \begin{tabular}{lll}
            \toprule
            \textbf{Week} & \textbf{Date} & \textbf{What Happens} \\
            \midrule
            9 & Apr. 24 & No class --- free research time \\
            10 & May 01 & No class --- free research time \\
            --- & \textbf{May 06} & \textbf{Assignment \#3 due} \\
            11 & May 08 & In-class review and consultations \\
            --- & \textbf{Jun. 01} & \textbf{Final manuscript due} \\
            \bottomrule
        \end{tabular}
    \end{table}
    \vspace{0.3cm}
    \textbf{What to focus on in the next three weeks:}
    \begin{enumerate}
        \item Finish data collection (if not done)
        \item Begin analysis --- apply your framework to your data
        \item Write the empirical chapter draft
        \item Meet with your supervisor
    \end{enumerate}
\end{frame}

\begin{frame}{Key Takeaways}
    \begin{enumerate}
        \item \textbf{Know your data.} If you cannot describe what you have, where it is, and how you will analyze it, fix that this week.
        \vspace{0.2cm}
        \item \textbf{Be transparent.} About your methods, your sources, and your use of any tools---including AI.
        \vspace{0.2cm}
        \item \textbf{When in doubt, ask.} Your supervisor is your first resource for questions about AI use, methodology, or scope.
        \vspace{0.2cm}
        \item \textbf{Assignment \#3 is the core of your thesis.} It shows you can do the research, not just plan it.
    \end{enumerate}
    \vspace{0.5cm}
    \begin{block}{Resources}
        \footnotesize
        \textcolor{SlateBlue}{\href{https://scdenney.github.io/baks_thesis-seminar/}{Course website}} ~|~
        \textcolor{SlateBlue}{\href{https://scdenney.github.io/baks_thesis-seminar/protocol/}{Thesis Protocol}} ~|~
        \textcolor{SlateBlue}{\href{https://www.organisatiegids.universiteitleiden.nl/en/regulations/humanities/guidelines-for-the-use-of-genai-in-assessment}{Faculty GenAI Guidelines}} ~|~
        \textcolor{SlateBlue}{\href{https://www.organisatiegids.universiteitleiden.nl/en/regulations/general/plagiarism}{Plagiarism Regulations}}
    \end{block}
\end{frame}

\end{document}
